% Options for packages loaded elsewhere
% Options for packages loaded elsewhere
\PassOptionsToPackage{unicode}{hyperref}
\PassOptionsToPackage{hyphens}{url}
\PassOptionsToPackage{dvipsnames,svgnames,x11names}{xcolor}
%
\documentclass[
  10pt,
  letterpaper,
  DIV=11,
  numbers=noendperiod]{scrartcl}
\usepackage{xcolor}
\usepackage{amsmath,amssymb}
\setcounter{secnumdepth}{5}
\usepackage{iftex}
\ifPDFTeX
  \usepackage[T1]{fontenc}
  \usepackage[utf8]{inputenc}
  \usepackage{textcomp} % provide euro and other symbols
\else % if luatex or xetex
  \usepackage{unicode-math} % this also loads fontspec
  \defaultfontfeatures{Scale=MatchLowercase}
  \defaultfontfeatures[\rmfamily]{Ligatures=TeX,Scale=1}
\fi
\usepackage{lmodern}
\ifPDFTeX\else
  % xetex/luatex font selection
\fi
% Use upquote if available, for straight quotes in verbatim environments
\IfFileExists{upquote.sty}{\usepackage{upquote}}{}
\IfFileExists{microtype.sty}{% use microtype if available
  \usepackage[]{microtype}
  \UseMicrotypeSet[protrusion]{basicmath} % disable protrusion for tt fonts
}{}
% Make \paragraph and \subparagraph free-standing
\makeatletter
\ifx\paragraph\undefined\else
  \let\oldparagraph\paragraph
  \renewcommand{\paragraph}{
    \@ifstar
      \xxxParagraphStar
      \xxxParagraphNoStar
  }
  \newcommand{\xxxParagraphStar}[1]{\oldparagraph*{#1}\mbox{}}
  \newcommand{\xxxParagraphNoStar}[1]{\oldparagraph{#1}\mbox{}}
\fi
\ifx\subparagraph\undefined\else
  \let\oldsubparagraph\subparagraph
  \renewcommand{\subparagraph}{
    \@ifstar
      \xxxSubParagraphStar
      \xxxSubParagraphNoStar
  }
  \newcommand{\xxxSubParagraphStar}[1]{\oldsubparagraph*{#1}\mbox{}}
  \newcommand{\xxxSubParagraphNoStar}[1]{\oldsubparagraph{#1}\mbox{}}
\fi
\makeatother


\usepackage{longtable,booktabs,array}
\usepackage{calc} % for calculating minipage widths
% Correct order of tables after \paragraph or \subparagraph
\usepackage{etoolbox}
\makeatletter
\patchcmd\longtable{\par}{\if@noskipsec\mbox{}\fi\par}{}{}
\makeatother
% Allow footnotes in longtable head/foot
\IfFileExists{footnotehyper.sty}{\usepackage{footnotehyper}}{\usepackage{footnote}}
\makesavenoteenv{longtable}
\usepackage{graphicx}
\makeatletter
\newsavebox\pandoc@box
\newcommand*\pandocbounded[1]{% scales image to fit in text height/width
  \sbox\pandoc@box{#1}%
  \Gscale@div\@tempa{\textheight}{\dimexpr\ht\pandoc@box+\dp\pandoc@box\relax}%
  \Gscale@div\@tempb{\linewidth}{\wd\pandoc@box}%
  \ifdim\@tempb\p@<\@tempa\p@\let\@tempa\@tempb\fi% select the smaller of both
  \ifdim\@tempa\p@<\p@\scalebox{\@tempa}{\usebox\pandoc@box}%
  \else\usebox{\pandoc@box}%
  \fi%
}
% Set default figure placement to htbp
\def\fps@figure{htbp}
\makeatother





\setlength{\emergencystretch}{3em} % prevent overfull lines

\providecommand{\tightlist}{%
  \setlength{\itemsep}{0pt}\setlength{\parskip}{0pt}}



 


\usepackage{xcolor, soul}
\usepackage{bm}
\definecolor{Light}{gray}{.90}
\sethlcolor{Light}
\let\OldTexttt\texttt
\renewcommand{\texttt}[1]{\OldTexttt{\hl{#1}}}
\usepackage{caption}
\captionsetup{font=small, labelfont=bf}
\setlength{\parindent}{0pt}
\setlength{\parskip}{1.25ex}
\newcommand{\bfX}{\mathbf{X}}
\newcommand{\bfY}{\mathbf{Y}}
\newcommand{\bfx}{\mathbf{x}}
\newcommand{\bfy}{\mathbf{y}}
\DeclareMathOperator{\diag}{diag}
\KOMAoption{captions}{tableheading}
\makeatletter
\@ifpackageloaded{caption}{}{\usepackage{caption}}
\AtBeginDocument{%
\ifdefined\contentsname
  \renewcommand*\contentsname{Table of contents}
\else
  \newcommand\contentsname{Table of contents}
\fi
\ifdefined\listfigurename
  \renewcommand*\listfigurename{List of Figures}
\else
  \newcommand\listfigurename{List of Figures}
\fi
\ifdefined\listtablename
  \renewcommand*\listtablename{List of Tables}
\else
  \newcommand\listtablename{List of Tables}
\fi
\ifdefined\figurename
  \renewcommand*\figurename{Figure}
\else
  \newcommand\figurename{Figure}
\fi
\ifdefined\tablename
  \renewcommand*\tablename{Table}
\else
  \newcommand\tablename{Table}
\fi
}
\@ifpackageloaded{float}{}{\usepackage{float}}
\floatstyle{ruled}
\@ifundefined{c@chapter}{\newfloat{codelisting}{h}{lop}}{\newfloat{codelisting}{h}{lop}[chapter]}
\floatname{codelisting}{Listing}
\newcommand*\listoflistings{\listof{codelisting}{List of Listings}}
\makeatother
\makeatletter
\makeatother
\makeatletter
\@ifpackageloaded{caption}{}{\usepackage{caption}}
\@ifpackageloaded{subcaption}{}{\usepackage{subcaption}}
\makeatother
\usepackage{bookmark}
\IfFileExists{xurl.sty}{\usepackage{xurl}}{} % add URL line breaks if available
\urlstyle{same}
\hypersetup{
  pdftitle={Normalization of neural responses},
  pdfauthor={Al Xin},
  colorlinks=true,
  linkcolor={blue},
  filecolor={Maroon},
  citecolor={Blue},
  urlcolor={Blue},
  pdfcreator={LaTeX via pandoc}}


\title{Normalization of neural responses}
\author{\small Al Xin}
\date{Aug 20, 2026}
\begin{document}
\maketitle


\section{Methods}\label{methods}

\subsection{Notation}\label{notation}

Let:

\begin{itemize}
\tightlist
\item
  \(n\): the number of sample stimuli.

  \begin{itemize}
  \tightlist
  \item
    \(n_f\): the number of texture families. Here, assume noisy texture
    families are omitted, so \(n_f = 35\)
  \item
    \(n_s\): the number of samples within each texture family.
  \item
    \(n_r\): the number of repetitions of each specific sample.
  \item
    \(n_b\): the number of blank stimuli.
  \end{itemize}
\item
  \(k_i\): the number of recording sites for region \(i\).

  \begin{itemize}
  \tightlist
  \item
    If needed, this can also be sectioned by array/simultaneously
    recorded sites.
  \end{itemize}
\item
  \(t\): the number of 10 ms bins of multi-unit activation (MUA) counts.
\end{itemize}

The raw MUAs for a given region are contained in a tensor of size
\(n \times k_i \times t\). For most analyses, we will summarize over
\(t\) to produce a two-dimensional matrix \(\bfX\). Then, \(X_{i, j}\)
is a specific element, \(\bfX_{-, i}\) is a \(n \times 1\) column vector
(the responses of one site to all stimuli), and \(\bfX_{i, -}\) is a
\(1 \times k_i\) row vector (the responses elicited across all sites by
one stimulus).

\subsection{Mahalanobis distance}\label{mahalanobis-distance}

Let \(P\) be some probability distributions on \(\mathbb{R}^k\) with
mean \(\pmb{\mu}\) and covariance matrix \(\pmb\Sigma\). The Mahalanobis
distance of points \(\bfx, \bfy\) with respect to \(P\) is \[
d_M (\bfx, \bfy, P) = 
\sqrt{ ( \bfx - \bfy )^\intercal \pmb\Sigma^{-1} ( \bfx - \bfy ) }.
\]

If \(\pmb\Sigma\) is positive semi-definite, the covariance can be
decomposed as \(\pmb\Sigma = \mathbf{S}^\intercal \mathbf{S}\), where
\(\mathbf{S}\) is the standard deviation matrix. So, the Mahalanobis
distance is equivalently \[
d_M (\bfx, \bfy, P) = \| \mathbf{S}^{-1} (\bfx - \bfy) \|, 
\] or the Euclidean distance after a whitening transformation.

\subsubsection{\texorpdfstring{Usage in Hotelling's two-sample
\emph{t}-test}{Usage in Hotelling's two-sample t-test}}\label{usage-in-hotellings-two-sample-t-test}

Let \(P \sim \mathcal{N}_k (\pmb\mu, \pmb\Sigma)\), i.e., a
\(k\)-dimensional multivariate Normal distribution. Then, let
\(\bfx_1, \bfx_2, \dots, \bfx_{n_x}\) and
\(\bfy_1, \bfy_2, \dots, \bfy_{n_y}\) be independent samples from \(P\).
Let each sample be a column vector. We can then define the sample mean
vectors \[
\overline{\bfx} = \frac{1}{n_x} \sum_{i = 1}^{n_x} \bfx_i, \quad
\overline{\bfy} = \frac{1}{n_y} \sum_{i = 1}^{n_y} \bfy_i,
\] and the sample covariance matrices as \[
\hat{\pmb\Sigma}_{\bfx} = \frac{1}{n_x - 1} \sum_{i = 1}^{n_x} (\bfx_i - \overline\bfx) (\bfx_i - \overline\bfx)^\intercal, \quad
\hat{\pmb\Sigma}_{\bfy} = \frac{1}{n_y - 1} \sum_{i = 1}^{n_y} (\bfy_i - \overline\bfy) (\bfy_i - \overline\bfy)^\intercal. 
\] The unbiased pooled covariance matrix estimate is the weighted
average \[
\hat{\pmb\Sigma} = 
\frac{ (n_x - 1) \hat{\pmb\Sigma}_{\bfx} + (n_y - 1) \hat{\pmb\Sigma}_{\bfy} }
{ n_x + n_y - 2 },
\] which is simply
\(\left( \hat{\pmb\Sigma}_{\bfx} + \hat{\pmb\Sigma}_{\bfy} \right) / 2\)
when \(n_x = n_y\).

Though not used in our analyses, the resulting Hotelling's two-sample
\emph{t}-squared statistic has the distribution \[
t^2 = \frac{n_x n_y}{n_x + n_y} (\overline{\bfx} - \overline{\bfy})^\intercal \hat{\pmb\Sigma}^{-1} (\overline{\bfx} - \overline{\bfy}) \sim T^2 (k, n_x + n_y - 2).
\]

\subsubsection{Effect of normalization on Mahalanobis distance
calculations}\label{effect-of-normalization-on-mahalanobis-distance-calculations}

As mentioned previously, the Mahalanobis distance is equivalent to the
Euclidean distance after whitening with \({\pmb{\Sigma}}^{-1/2}\). Thus,
any data processing that can be ``undone'' by whitening will not change
the outcome. For instance, centering and scaling with respect to the
sites will result in the same value.

Let \(\mathbf{b} \in \mathbb{R}^k\) be some arbitrary constant vector
and \(\mathbf{m} \in \mathbb{R}^k\) be some non-zero constant vector. If
we take \[
\begin{aligned}
\bfx' &= \begin{bmatrix} m_1 x_1 + b_1 & m_2 x_2 + b_2 & \cdots & m_k x_k + b_k  \end{bmatrix}^\intercal, \\
\bfy' &= \begin{bmatrix} m_1 y_1 + b_1 & m_2 y_2 + b_2 & \cdots & m_k y_k + b_k  \end{bmatrix}^\intercal,
\end{aligned}
\] then transform the assumed covariance matrix of the new distribution
\(P\) accordingly with \[
\tilde{\pmb{\Sigma}} = \diag (\mathbf{m}) \pmb\Sigma \diag (\mathbf{m}),
\] then the term in the Mahalanobis distance simplifies to \[
\begin{aligned}
(\bfx' - \bfy')^\intercal \tilde{\pmb\Sigma}^{-1} (\bfx' - \bfy') &= 
(\bfx - \bfy)^\intercal \diag (\mathbf{m}) \left(  \diag (\mathbf{m}) \pmb\Sigma \diag (\mathbf{m}) \right)^{-1} \diag (\mathbf{m}) (\bfx - \bfy) \\
&= (\bfx - \bfy)^\intercal \pmb\Sigma^{-1} (\bfx - \bfy).
\end{aligned}
\]

Though the sample setting estimates the covariance rather than relying
on the known \(P\), this invariance to centering and scaling still
holds. That is,

\[
\begin{aligned}
\hat{\pmb\Sigma}_{\bfx'} &= \frac{1}{n_x - 1} \sum_{i = 1}^{n_x} (\bfx_i' - \overline{\bfx}') (\bfx_i' - \overline{\bfx}')^\intercal \\
&= \frac{1}{n_x - 1} \sum_{i = 1}^{n_x} \diag(\mathbf{m}) (\bfx_i - \overline{\bfx}) (\bfx_i - \overline{\bfx})^\intercal \diag (\mathbf{m}) \\
\bfx' - \bfy' &= \diag (\mathbf{m}) (\bfx - \bfy)
\end{aligned}
\]

The effect of centering and scaling on \(\hat\Sigma\) can still be
reduced to multiplication by a diagonal matrix.

\subsubsection{Pairwise between-family
distances}\label{pairwise-between-family-distances}

When calculating the distance between each pair of texture families,
each sample image and each repetition can be considered a draw from the
underlying distribution.




\end{document}
